\documentclass[10pt,a4paper]{amsart}
\usepackage[margin=19mm]{geometry}
\usepackage{amsmath,amssymb,mathtools}
\usepackage[scr=rsfs,cal=euler]{mathalfa}
\usepackage{xfrac}
\usepackage[unicode,hidelinks,bookmarks=false]{hyperref}
\newcommand{\ie}{i.e.\ }
\newcommand{\cf}{cf.\ }
\newcommand{\resp}{respectively\ }
\newcommand{\Subsection}{\subsection}
\providecommand{\nobreakdash}{\nobreak\hbox{-}}
\setlength{\emergencystretch}{2em}
\multlinegap0pt
\usepackage{stackengine}
\usepackage{amssymb}
\usepackage{mathtools}
\usepackage{tikz}
\usepackage[shortlabels]{enumitem}
\usepackage{csquotes}
\newtheorem{prop}{Proposition}
\newtheorem{thm}{Theorem}
\usepackage{cleveref}
\crefname{prop}{Proposition}{Propositions}
\crefname{thm}{Theorem}{Theorems}
\renewcommand{\L}{\mathcal{L}}
\newcommand{\LWitt}{\L_{\mathrm{Witt}}}
\newcommand{\Loag}{\L_{\mathrm{oag}}}
\newcommand{\Lring}{\L_{\mathrm{ring}}}
\newcommand{\Lval}{\L_{\mathrm{val}}}
\renewcommand{\O}{\mathcal{O}}
\DeclareMathOperator{\res}{res}
\title{AKE principles in roughly deeply ramified henselian valued fields}
\author{Franziska Jahnke, Margarete Ketelsen, Floris Vermeulen}
\date{Source: arXiv:2609.15376v1 (CC BY 4.0)}
\begin{document}
\maketitle

\noindent\emph{Source and licence.} AKE principles in roughly deeply ramified henselian valued fields. Version arXiv:2609.15376v1 (CC BY 4.0), distributed by arXiv under the Creative Commons Attribution 4.0 International licence, http://creativecommons.org/licenses/by/4.0/, as recorded in the arXiv licence metadata for this exact version and shown on the abstract page https://arxiv.org/abs/2609.15376v1. Full licence notice: \texttt{licenses/arxiv-2609.15376-CC-BY-4.0.txt}.\par\smallskip

\noindent\textit{Editorial scope.} Counted unit: the article's thm thm:main (source0.tex lines 698--709). The article's complete original proof of it is delivered with it (source0.tex lines 711--786, one \texttt{proof} environment of 76 lines). No mathematics is written, repaired or re-derived: the statement and the proof are the article's own lines, in the article's own notation, and no retained line is substituted or re-flowed.  Retained uncounted as the source-local context the proof needs: lines 620--627.  Nothing was omitted from any retained span, so the package carries no omission record.  The retained lines cite 4 works, whose entries are transcribed verbatim from the article's own bibliography source in the e-print and delivered with short editorial labels recorded in the item's reference mapping.  No retained line contains a figure environment, an image-inclusion command or drawing code, so no image is delivered and the package declares no asset dependency.  Editorial formatting only, in addition: an AMS \texttt{amsart} preamble that restates the article's own declarations and macros used by the retained lines, namely the environments \texttt{prop}, \texttt{thm} on separate counters, together with the standard packages those lines need.  The retained environments print in this document as Proposition 1, Theorem 1 in order of appearance, whereas the article prints the same items with the numbering of its own full document.  The complete native source of the article as distributed in the arXiv e-print for this exact version is bundled unchanged as \texttt{sources/210411/source0.tex}.
    \begin{prop} \label{prop:lift}
        Let $(K,v)$ and $(L,w)$ be $p$-adically complete valued fields of mixed characteristic $(0,p)$
        such that $\O_v/(p)$ and $\O_w/(p)$ are semiperfect. 
        Then any embedding (respectively isomorphism)  
        $\O_v/(p) \to \O_w/(p)$ as $\LWitt$-structures is induced by a
        unique embedding (respectively
        isomorphism) $(K,v) \to (L,w).$
    \end{prop}

    \begin{thm}
        \label{thm:main}
        Let $(K,v)$ and $(L,w)$ be henselian valued fields of mixed characteristic 
        $(0,p)$ 
        such that $\O_v/(p)$ and $\O_w/(p)$ are semiperfect. Then, we have
        $$\underbrace{(K,v) \equiv (L,w)}_{\text{in }\mathcal{L}_\mathrm{val}} \Longleftrightarrow \underbrace{\mathcal{O}_v/(p) \equiv \mathcal{O}_w/(p)}_{\text{in }\mathcal{L}_{\mathrm{Witt}}}\textrm{ and } \underbrace{(vK,v(p)) 
        \equiv (wL,w(p))}_{\text{in }\mathcal{L}_{\mathrm{oag}}\cup\{c\}},$$
        and 
        $$\underbrace{(K,v) \equiv_\exists (L,w)}_{\text{in }\mathcal{L}_\mathrm{val}} \Longleftrightarrow \underbrace{\mathcal{O}_v/(p) \equiv_\exists \mathcal{O}_w/(p)}_{\text{in }\mathcal{L}_{\mathrm{Witt}}}\textrm{ and } \underbrace{(vK,v(p)) \equiv_\exists (wL,w(p))}_{\text{in }\mathcal{L}_{\mathrm{oag}}\cup\{c\}},$$
        where $c$ is a constant symbol interpreted as the value of $p$.
        Moreover, when $vK$ and $wL$ are regular, the value group data on the right hand side can be omitted in both equivalences.
    \end{thm}

    \begin{proof}
        If $(K,v)\equiv (L, w)$ then clearly we also have that $\O_v/(p) \equiv \O_w/(p)$ in $\LWitt$ and $(vK, v(p))\equiv (wL, w(p))$ in $\Loag\cup\{c\}$.
        Indeed, this follows from the fact that $\O_v/(p)$ and $(vK, v(p))$ with their respective languages are interpretable in $(K,v)$.
        If $\mathrm{Th}_\exists^{\Lval}{(K,v)} \subseteq \mathrm{Th}_\exists^{\Lval} (L,w)$, then -- potentially after replacing $(L,w)$ by an elementary extension -- we may assume that we are given an $\Lval$-embedding $\iota\colon K\to L$ (see \cite[Proposition 5.2.2]{CK}).
        It is clear that $\iota$ induces an $\Lring$-embedding $\O_v/(p)\to \O_w/(p)$ and an $\Loag\cup\{c\}$-embedding $(vK, v(p))\to (wL, w(p))$.
        We wish to show that the map $\O_v/(p)\to \O_w/(p)$ is in fact an $\LWitt$-embedding.
        For $x_0, \ldots, x_{n-1}, y_0, \ldots, y_{n-1}\in \O_v$ we have that
        \begin{align*}
        \ker_n(x_0+& p\O_v, \ldots, x_{n-1}+p\O_v; y_0+p\O_v, \ldots, y_{n-1}+p\O_v) \\
        &\Longleftrightarrow v\big(W_{n-1}(x_0, \ldots, x_{n-1}) - W_{n-1}(y_0, \ldots, y_{n-1})\big) \geq v(p^n).    
        \end{align*}
        Clearly this condition is preserved by $\iota$.
        Thus, we conclude $$\mathrm{Th}_\exists^{\LWitt}(\O_v/(p))\subseteq \mathrm{Th}_{\exists}^{\LWitt}(\O_w/(p)).$$
        By symmetry, we obtain
        $$\underbrace{(K,v) \equiv_\exists (L,w)}_{\text{in }\mathcal{L}_\mathrm{val}} \Longrightarrow \underbrace{\mathcal{O}_v/(p) \equiv_\exists \mathcal{O}_w/(p)}_{\text{in }\mathcal{L}_{\mathrm{Witt}}}\textrm{ and } \underbrace{(vK,v(p)) \equiv_\exists (wL,w(p))}_{\text{in }\mathcal{L}_{\mathrm{oag}}\cup\{c\}}.$$
        
        We now focus on the other direction. We assume that 
        $$
        \underbrace{\mathcal{O}_v/(p) \equiv_{(\exists)} \mathcal{O}_w/(p)}_{\text{in }\mathcal{L}_{\mathrm{Witt}}}\textrm{ and } \underbrace{(vK,v(p)) \equiv_{(\exists)} (wL,w(p))}_{\text{in }\mathcal{L}_{\mathrm{oag}}\cup\{c\}}.$$
        holds and -- without loss of generality -- further 
        that $(K,v)$ and $(L,w)$ are $\aleph_1$-saturated.
        In the $\equiv_\exists$-case, invoking \cite[Proposition 5.2.2]{CK} once again,
        we may assume there are embeddings
        \[
        f\colon \O_v/(p)\to \O_w/(p), \quad g\colon (vK, v(p))\to (wL, w(p))
        \]
        in the language $\LWitt$ resp.\ $\Loag\cup\{c\}$.
        In the $\equiv$-case, by the Keisler--Shelah theorem \cite[Theorem 6.1.15]{CK}, we may moreover assume that both
        $f$ and $g$ are isomorphisms.

        Let $\Delta_K\leq vK$ be the smallest convex subgroup of $vK$ containing $v(p)$, and similarly define $\Delta_L\leq wL$.
        Let $v_0\colon K^\times\to vK / \Delta_K$ and $w_0\colon L^\times\to wL / \Delta_L$ be the corresponding coarsenings of the valuations.
        By saturation, the residue field of $v_0$ is 
        \[K_0\cong \mathrm{Frac}(\varprojlim \O_v/(p^n))\]
        with induced valuation $\overline{v}\colon K_0^\times\to \Delta_K$, and similarly for the residue field $L_0$ of $w_0$.
        In other words, $K_0$ and $L_0$ are $p$-adically complete and both $\O_{\overline{v}}/(p) = \O_v/(p)$ and $\O_{\overline{w}}/(p) = \O_w/(p)$
         are  semiperfect.
        Hence by \Cref{prop:lift}, the map $f\colon \O_v/(p)\to \O_w/(p)$ lifts to give an isomorphism (resp.\ embedding)
        \[
        f'\colon (K_0, \overline{v})\to (L_0, \overline{w}).
        \]

        Since $g\colon (vK, v(p))\to (wL, w(p))$ maps $v(p)$ to $w(p)$, it descends to an isomorphism (resp.\ embedding)
        \[
        g'\colon vK/\Delta_K\to wL /\Delta_L. 
        \]

        We now make a case distinction:
        \begin{enumerate}
        \item In case both $f'$ and $g'$ are isomorphisms, the 
        Ax--Kochen--Ershov theorem in equicharacteristic zero shows 
        \[
         (K,v_0)\equiv (L, w_0).
        \]
        By resplendency (see \cite[Proposition 2.3]{ketelsenCompositionAxKochenErshov2026}), this implies 
        $(K,v)\equiv (L,w)$, as desired.
        
        \item In case $f'$ or $g'$ is not an isomorphism (but only an embedding), consider an $|K|^+$-saturated elementary extension $(M,\nu, \omega)$
        of the bi-valued field $(L,w_0,w)$. Then $f'$ and $g'$ give rise to embeddings
        $$f'' \colon (K_0,\overline{v}) \hookrightarrow (M\nu, \overline{\omega})
        \textrm{ and } g'' \colon vK/\Delta_K = v_0K \hookrightarrow \nu M$$
        where $\overline{\omega}$ is the valuation induced by \(\omega\) on $M\nu$.
        By the Embedding Lemma for henselian valued fields of equicharacteristic $0$, see \cite[Lemma 4.6.2]{PD}, we obtain an embedding $h: (K,v_0) \hookrightarrow (M,\nu)$ inducing both $f''$ and $g''$. As for any $x \in K$, we have
        \begin{align*}
        x \in \O_v &\Longleftrightarrow {\res}_{v_0}(x) \in \O_{\overline{v}} \Longleftrightarrow f'(\res_{w_0}(x)) \in \O_{\overline{w}} \\ & \Longleftrightarrow f''(\res_\nu (x)) \in \O_{\overline{\omega}} ,
        \end{align*}
        $h$ is automatically also an embedding of $(K,v)$ into $(M,\omega) \succeq (L,w)$. In particular, $\mathrm{Th}_\exists(K,v) \subseteq \mathrm{Th}_\exists(L,w)$.
        Symmetrically, starting with appropriate 
        embeddings of $\O_w/(p) \hookrightarrow \O_v/(p)$
        and $(wL,w(p)) \hookrightarrow (vK,v(p))$ we obtain $\mathrm{Th}_\exists(K,v) \supseteq \mathrm{Th}_\exists(L,w)$.
        \end{enumerate}

        For the moreover statement: $\aleph_1$-saturation shows that $vK/\Delta_K$ and $wL/\Delta_L$ are nontrivial.
        As regularity implies that $vK/\Delta_K$ and $wL/\Delta_L$ are divisible, see
        \cite[\S 3]{Rob}, we automatically have that $vK/\Delta_K \equiv_{(\exists)} wL / \Delta_L$.
    \end{proof}

\begin{thebibliography}{99}
\bibitem[CK]{CK} Chen~Chung Chang and Howard~Jerome Keisler, \emph{Model theory}, third ed., Studies in Logic and the Foundations of Mathematics, vol.~73, North-Holland, Amsterdam and New York, 1990.
\bibitem[PD]{PD} Alexander Prestel and Charles~N. Delzell, \emph{Mathematical logic and model theory}, Universitext, Springer, London, 2011.
\bibitem[Rob]{Rob} Abraham Robinson and Elias Zakon, \emph{Elementary properties of ordered abelian groups}, Trans. Amer. Math. Soc. \textbf{96} (1960), 222--236.
\bibitem[ketels]{ketelsenCompositionAxKochenErshov2026} Margarete Ketelsen, \emph{Composition {A}x--{K}ochen/{E}rshov principles and tame fields of mixed characteristic}, 2026, Preprint, arXiv:2601.05790, with an appendix by Philip Dittmann.
\end{thebibliography}
\end{document}
